Long-horizon and continual-learning agents keep what they learn in memory that a language model
repeatedly rewrites: after each session the model reads its previous memory and writes a condensed
replacement, and the original text is gone. Treating this as an iterated-learning chain predicts
that survival should track the compactor's prior, so memory would be worst at keeping exactly the
content it exists for --- a claim repeated informally across the literature and in production
documentation, but never tested with prior-fit as a manipulated variable. We test it. Using
minimal-pair memory items whose prior-fit is graded and measured on the compactor itself by paired
comparison, and a same-item manipulation that holds the target sentence \emph{and its value} fixed
while changing only a value-free description of the surrounding work, we find \textbf{no support for
the predicted effect}: making the very same value expected rather than surprising moves survival by
\EShiftMean{}~pp (95\% CI [\EShiftCIlo{}, \EShiftCIhi{}]), ruling out benefits larger than about
\ConeUpper{}~pp for this manipulation. The graded dose--response is likewise consistent with zero
(\DoseIQR{}~pp across the interquartile surprise contrast, 95\% CI [\DoseIQRlo{}, \DoseIQRhi{}]), and
every slope interval in our portability diagnostics --- a second strong model family, a small local
model, and temperature $0$ --- contains zero. Nor is memory corrupted: across \NCodedCases{} coded
item-checkpoints only \ChangedCount{} are changes to a different value; under a binding budget the
compactor simply deletes. What separates operators is much larger than anything surprise explains.
Rewriting memory by reconstruction from the retained raw sessions keeps \RetOfour{}\% of items after
ten compactions, and adding the anti-drift rules a production memory system already ships keeps
\RetOfive{}\%, against \RetOone{}\% for recursive self-rewriting at the same budget
(\ContrastOfour{} and \ContrastOfive{}~pp). Pinning high-surprise items --- the fix our own theory
predicted --- shows no detected benefit and is indistinguishable from token-matched random pinning.
Preferences show no homogenisation either: those that differ from the model's defaults survive
\PrefAdj{}~pp better after baseline adjustment (95\% CI [\PrefAdjLo{}, \PrefAdjHi{}]), so stored
profiles become more distinctive. We release the instrument, the stimuli and all result files, and we
report the ways our own causal manipulation is imperfect.
